\begin{derivation}[从\eqnrefs{eq:JK-algebra}到\eqnrefs{eq:lorentz-su2-algebra-dp34}]

正文\eqnrefs{eq:JK-algebra} 已把六个 Lorentz 生成元写成三维转动 $J_i$ 与推动 $K_i$。将它们组合成 $\mathsf J_i^{(\pm)}=(J_i\pm\ii K_i)/2$ 后，逐项代入三组 $J$--$K$ 对易关系即可证明两组生成元各自闭合且彼此对易，从而得到正文\eqnrefs{eq:lorentz-su2-algebra-dp34}。



由
\begin{equation*}
 \mathsf J_i^{(+)}=\frac12(J_i+\ii K_i),
 \qquad
 \mathsf J_i^{(-)}=\frac12(J_i-\ii K_i)
\end{equation*}
与
\begin{equation*}
 [J_i,J_j]=\ii\hbar\epsilon_{ijk}J_k,
 \qquad
 [J_i,K_j]=\ii\hbar\epsilon_{ijk}K_k,
 \qquad
 [K_i,K_j]=-\ii\hbar\epsilon_{ijk}J_k,
\end{equation*}
先计算第一组：
\begin{align*}
 [\mathsf J_i^{(+)},\mathsf J_j^{(+)}]
 &=\frac14\Bigl(
 [J_i,J_j]+\ii[J_i,K_j]
 +\ii[K_i,J_j]-[K_i,K_j]\Bigr).
\end{align*}
利用
\begin{equation*}
 [K_i,J_j]=-[J_j,K_i]
 =\ii\hbar\epsilon_{ijk}K_k,
\end{equation*}
得到
\begin{align*}
 [\mathsf J_i^{(+)},\mathsf J_j^{(+)}]
 &=\frac14\left(
 2\ii\hbar\epsilon_{ijk}J_k
 -2\hbar\epsilon_{ijk}K_k\right),\\
 &=\ii\hbar\epsilon_{ijk}
 \frac{J_k+\ii K_k}{2},\\
 &=\ii\hbar\epsilon_{ijk}\mathsf J_k^{(+)}.
\end{align*}
第二组从定义 $\mathsf J_i^{(-)}=(J_i-\ii K_i)/2$ 独立计算：
\begin{align*}
 [\mathsf J_i^{(-)},\mathsf J_j^{(-)}]
 &=\frac14\Bigl(
 [J_i,J_j]-\ii[J_i,K_j]
 -\ii[K_i,J_j]-[K_i,K_j]\Bigr),\\
 &=\frac14\Bigl(
 \ii\hbar\epsilon_{ijk}J_k
 +\hbar\epsilon_{ijk}K_k
 +\hbar\epsilon_{ijk}K_k
 +\ii\hbar\epsilon_{ijk}J_k\Bigr),\\
 &=\frac12\Bigl(
 \ii\hbar\epsilon_{ijk}J_k
 +\hbar\epsilon_{ijk}K_k\Bigr),\\
 &=\ii\hbar\epsilon_{ijk}\frac{J_k-\ii K_k}{2},\\
 &=\ii\hbar\epsilon_{ijk}\mathsf J_k^{(-)}.
\end{align*}
交叉对易子为
\begin{align*}
 [\mathsf J_i^{(+)},\mathsf J_j^{(-)}]
 &=\frac14\Bigl(
 [J_i,J_j]-\ii[J_i,K_j]
 +\ii[K_i,J_j]+[K_i,K_j]\Bigr),\\
 &=0.
\end{align*}
因此复化 Lorentz 李代数确实分解为两份彼此对易的 $\mathfrak{su}(2)$。
\end{derivation}

\begin{derivation}[从\eqnrefs{eq:sigmabar-identity}到\eqnrefs{eq:massive-weyl-coupled-r7}]

正文已指出 $p_\mu\sigma^\mu$ 与 $p_\mu\bar\sigma^\mu$ 是两个手征表示之间的一阶 Lorentz 交织映射。这里只需确定允许的标量系数，并由质量壳固定它们的乘积；静止系不参与一般式的建立。

先写最一般线性耦合
\begin{equation*}
 (p_\mu\sigma^\mu)\psi_R=\mu_L\psi_L,
 \qquad
 (p_\mu\bar\sigma^\mu)\psi_L=\mu_R\psi_R.
\end{equation*}
由正文\eqnrefs{eq:sigmabar-identity}，
\begin{align*}
 (p_\mu\bar\sigma^\mu)(p_\nu\sigma^\nu)
 &=\frac12p_\mu p_\nu
 \left(\bar\sigma^\mu\sigma^\nu+\bar\sigma^\nu\sigma^\mu\right)\\
 &=p_\mu p^\mu\,\id_2.
\end{align*}
把第一式左乘 $p_\mu\bar\sigma^\mu$，再代入第二式：
\begin{align*}
 p^2\psi_R
 &=(p\!\cdot\!\bar\sigma)(p\!\cdot\!\sigma)\psi_R\\
 &=\mu_L(p\!\cdot\!\bar\sigma)\psi_L\\
 &=\mu_L\mu_R\psi_R.
\end{align*}
对非零态得到
\begin{equation*}
 p^2=\mu_L\mu_R.
\end{equation*}
同样作用在 $\psi_L$ 上得到相同条件。对质量壳 $p^2=\me^2c^2$，
\begin{equation*}
 \mu_L\mu_R=\me^2c^2.
\end{equation*}

若自由二次作用量厄米，左右两条方程中的质量系数互为复共轭，
\begin{equation*}
 \mu_R=\mu_L^*.
\end{equation*}
于是 $|\mu_L|=|\mu_R|=\me c$。对两个手征场作相反的常数重相位可消去这一个共同复相位，因此可选
\begin{equation*}
 \mu_L=\mu_R=\me c>0.
\end{equation*}
这就得到正文\eqnrefs{eq:massive-weyl-coupled-r7}。

作为随后才做的检查，把质量壳写成正文\eqnrefs{eq:rapidity-momentum-r7}。由 Pauli 代数，
\begin{align*}
 p_\mu\sigma^\mu
 &=\me c\left(\cosh\zeta\,\id_2
 -\sinh\zeta\,\hat{\mathbf n}\cdot\boldsymbol\sigma_{\rm P}\right)\\
 &=\me c\,\ee^{-\zeta\hat{\mathbf n}\cdot\boldsymbol\sigma_{\rm P}},\\
 p_\mu\bar\sigma^\mu
 &=\me c\,\ee^{+\zeta\hat{\mathbf n}\cdot\boldsymbol\sigma_{\rm P}}.
\end{align*}
与正文\eqnrefs{eq:weyl-boost-r7} 的左右手有限推动一致。快度表示在这里提供一般耦合式的一种显式质量壳参数化检验；一般方程本身已经由 Lorentz 交织映射与质量壳条件确定。
\end{derivation}



\begin{derivation}[从\eqnrefs{eq:clifford-general}到\eqnrefs{eq:dirac-minimal-spinor-dimension-dp68}]
设 $A,B$ 是任意两个不同的 Clifford 生成元，因此

\begin{equation*}
 AB=-BA,
 \qquad
 \det A\neq0,
 \qquad
 \det B\neq0.
\end{equation*}
若表示维数为奇数 $n$，则
\begin{align*}
 \det(AB)
 &=\det(-BA)\\
 &=(-1)^n\det(BA)\\
 &=-\det(AB),
\end{align*}
与 $\det(AB)\neq0$ 矛盾。因此任何有限维复 Clifford 表示的维数在这里都不能为奇数；特别地，三维表示已经被排除。

还需排除唯一更小的偶数候选 $n=2$。任意 $2\times2$ 复矩阵可唯一写成
\begin{equation*}
 X=a\id_2+\mathbf b\cdot\boldsymbol\sigma_{\rm P}.
\end{equation*}
若三个空间 Clifford 生成元存在于二维空间，经相似变换与非零常数重标度后可取为三条 Pauli 方向。与三者都反对易的时间生成元 $X$ 必须满足
\begin{equation*}
 0=\{X,\sigma_i\}
 =2a\sigma_i+2b_i\id_2
 \qquad(i=1,2,3).
\end{equation*}
取迹给 $b_i=0$；再乘任意 $\sigma_j$ 后取迹给 $a=0$。所以唯一同时与三个 Pauli 方向反对易的二维矩阵是零矩阵，这与 $(\gamma^0)^2=\id_2$ 矛盾。

因此非零复表示的维数既不能为奇数，也不能为二，最小可能值至少为四。正文\eqnrefs{eq:weyl-gamma} 又给出显式 $4\times4$ 矩阵并满足\eqnrefs{eq:clifford-general}，所以
\begin{equation*}
 \dim_{\mathbb C}S_{\rm Dirac}^{\rm min}=4.
\end{equation*}
这才完成“最小维数为四”的存在性与不可更小性两部分证明。
\end{derivation}



\begin{derivation}[从\eqnrefs{eq:sigma-definition}到\eqnrefs{eq:spinor-lorentz-algebra}]
先计算
\begin{align*}
 [\gamma^\mu\gamma^\nu,\gamma^\rho]
 &=\gamma^\mu\gamma^\nu\gamma^\rho
 -\gamma^\rho\gamma^\mu\gamma^\nu\\
 &=\gamma^\mu(2\eta^{\nu\rho}-\gamma^\rho\gamma^\nu)
 -(2\eta^{\mu\rho}-\gamma^\mu\gamma^\rho)\gamma^\nu\\
 &=2\eta^{\nu\rho}\gamma^\mu-2\eta^{\mu\rho}\gamma^\nu.
\end{align*}
交换 $\mu\leftrightarrow\nu$ 后相减，再乘正文\eqnrefs{eq:sigma-definition} 中的 $\ii/4$，便得到
\begin{equation*}
 [\Sigma^{\mu\nu},\gamma^\rho]
 =\ii(\eta^{\nu\rho}\gamma^\mu-\eta^{\mu\rho}\gamma^\nu),
\end{equation*}
即正文\eqnrefs{eq:sigma-gamma-comm}。

由 $\Sigma^{\rho\sigma}=\frac{\ii}{4}[\gamma^\rho,\gamma^\sigma]$，有
\begin{align*}
 [\Sigma^{\mu\nu},\Sigma^{\rho\sigma}]
 &=\frac{\ii}{4}\Bigl(
 [\Sigma^{\mu\nu},\gamma^\rho]\gamma^\sigma
 +\gamma^\rho[\Sigma^{\mu\nu},\gamma^\sigma]\\
 &\qquad-[\Sigma^{\mu\nu},\gamma^\sigma]\gamma^\rho
 -\gamma^\sigma[\Sigma^{\mu\nu},\gamma^\rho]\Bigr).
\end{align*}
代入正文\eqnrefs{eq:sigma-gamma-comm}，再把出现的 $\gamma^\alpha\gamma^\beta-\gamma^\beta\gamma^\alpha$ 重新写回 $\Sigma^{\alpha\beta}$，得到
\begin{equation*}
 [\Sigma^{\mu\nu},\Sigma^{\rho\sigma}]
 =\ii\left(
 \eta^{\nu\rho}\Sigma^{\mu\sigma}
 -\eta^{\mu\rho}\Sigma^{\nu\sigma}
 -\eta^{\nu\sigma}\Sigma^{\mu\rho}
 +\eta^{\mu\sigma}\Sigma^{\nu\rho}
 \right),
\end{equation*}
即正文\eqnrefs{eq:spinor-lorentz-algebra}。
\end{derivation}

\begin{derivation}[从\eqnrefs{eq:spinor-lorentz-transform}到\eqnrefs{eq:gamma-covariance}]

正文\eqnrefs{eq:spinor-lorentz-transform} 给出有限旋量变换的指数形式。其无穷小相似变换由 $[\Sigma^{\mu\nu},\gamma^\rho]$ 完全决定；正文\eqnrefs{eq:sigma-gamma-comm} 恰好把这一对易子写成四矢量表示的生成元作用，因此两种表示具有同一 Lie 代数微分，指数化后得到 $\gamma^\rho$ 的有限协变变换律。

写成一阶参数展开，
\begin{equation*}
 S=\id-\frac{\ii}{2}\omega_{\mu\nu}\Sigma^{\mu\nu}+O(\omega^2),
 \qquad
 S^{-1}=\id+\frac{\ii}{2}\omega_{\mu\nu}\Sigma^{\mu\nu}+O(\omega^2).
\end{equation*}
于是
\begin{align*}
 S^{-1}\gamma^\rho S
 &=\gamma^\rho+\frac{\ii}{2}\omega_{\mu\nu}
 [\Sigma^{\mu\nu},\gamma^\rho]+O(\omega^2)\\
 &=\gamma^\rho+\omega^\rho{}_{\sigma}\gamma^\sigma+O(\omega^2).
\end{align*}
四矢量表示同样满足
\begin{equation*}
 \Lambda^\rho{}_{\sigma}
 =\delta^\rho{}_{\sigma}+\omega^\rho{}_{\sigma}+O(\omega^2).
\end{equation*}
两种表示由同一 Lie 代数参数生成，因此沿群参数连续积分后得到
\begin{equation*}
 S^{-1}(\Lambda)\gamma^\rho S(\Lambda)
 =\Lambda^\rho{}_{\sigma}\gamma^\sigma,
\end{equation*}
即正文\eqnrefs{eq:gamma-covariance}。
\end{derivation}

\begin{derivation}[从\eqnrefs{eq:dirac-adjoint}到\eqnrefs{eq:dirac-adjoint-covariance-dp68}]

普通内积 $\psi^\dagger\psi$ 在推动下并不保持不变，因此正文\eqnrefs{eq:dirac-adjoint} 引入 $\gamma^0$ 来构造适合 Lorentz 群的双线性。要验证这个定义确实有效，只需证明 $S^\dagger\gamma^0=\gamma^0S^{-1}$。

在所采用的度规约定下，
\begin{equation*}
 (\gamma^0)^\dagger=\gamma^0,
 \qquad
 (\gamma^i)^\dagger=-\gamma^i,
\end{equation*}
等价于
\begin{equation*}
 (\gamma^\mu)^\dagger=\gamma^0\gamma^\mu\gamma^0.
\end{equation*}
因此
\begin{equation*}
 (\Sigma^{\mu\nu})^\dagger=\gamma^0\Sigma^{\mu\nu}\gamma^0.
\end{equation*}
对正文\eqnrefs{eq:spinor-lorentz-transform} 取厄米共轭便有
\begin{align*}
 S^\dagger
 &=\exp\!\left[+\frac{\ii}{2}\omega_{\mu\nu}(\Sigma^{\mu\nu})^\dagger\right]\\
 &=\gamma^0\exp\!\left[+\frac{\ii}{2}\omega_{\mu\nu}\Sigma^{\mu\nu}\right]\gamma^0
 =\gamma^0S^{-1}\gamma^0.
\end{align*}
于是
\begin{align*}
 \bar\psi'
 &=(S\psi)^\dagger\gamma^0
 =\psi^\dagger S^\dagger\gamma^0
 =\psi^\dagger\gamma^0S^{-1}
 =\bar\psi S^{-1},
\end{align*}
这就是正文\eqnrefs{eq:dirac-adjoint-covariance-dp68}。
\end{derivation}

\begin{derivation}[从\eqnrefs{eq:gamma5-def}到\eqnrefs{eq:chiral-projector-algebra-dp68}]
以 $\mu=0$ 为例，
\begin{align*}
 \gamma^5\gamma^0
 &=\ii\gamma^0\gamma^1\gamma^2\gamma^3\gamma^0\\
 &=-\ii\gamma^0\gamma^0\gamma^1\gamma^2\gamma^3
 =-\gamma^0\gamma^5.
\end{align*}
同理对任意 $\mu$ 都有 $\{\gamma^5,\gamma^\mu\}=0$。平方时，右半组四个 $\gamma$ 移到左半组同名矩阵旁共需 $3+2+1=6$ 次交换：
\begin{align*}
 (\gamma^5)^2
 &=-\gamma^0\gamma^1\gamma^2\gamma^3
 \gamma^0\gamma^1\gamma^2\gamma^3\\
 &=-(-1)^6(\gamma^0)^2(\gamma^1)^2(\gamma^2)^2(\gamma^3)^2
 =\id_4.
\end{align*}
两条关系合起来就是正文\eqnrefs{eq:gamma5-anticomm}。

先计算
\begin{align*}
 P_R^2&=\frac14(1+2\gamma^5+(\gamma^5)^2)=P_R,\\
 P_L^2&=\frac14(1-2\gamma^5+(\gamma^5)^2)=P_L,\\
 P_RP_L&=\frac14(1-(\gamma^5)^2)=0,\\
 P_R+P_L&=\id_4.
\end{align*}
再用 $\gamma^5\gamma^\mu=-\gamma^\mu\gamma^5$，
\begin{align*}
 P_R\gamma^\mu
 &=\frac12(1+\gamma^5)\gamma^\mu
 =\frac12\gamma^\mu(1-\gamma^5)
 =\gamma^\mu P_L,\\
 P_L\gamma^\mu&=\gamma^\mu P_R.
\end{align*}
这就得到正文\eqnrefs{eq:chiral-projector-algebra-dp68}。
\end{derivation}
